% reality/part_intro.tex

\chapter*{Reality}
\phantomsection
\addcontentsline{toc}{chapter}{Reality}

% ============================================================
\section*{What do we already know?}
% ============================================================

\begin{remark}
The previous parts established a complete mathematical description of
railway alignment geometry, railway states, runtime interpretation, and
operational evaluation.
\end{remark}

\begin{remark}
Within AIM, geometry emerges from intrinsic curvature evolution,
railway states describe spatial configuration, and runtime services
generate operational quantities through evaluation operators.
\end{remark}

\begin{remark}
These constructions provide a coherent idealized railway model.
\end{remark}

% ============================================================
\section*{What is still missing?}
% ============================================================

\begin{remark}
Real railway systems do not exist purely within mathematical models.
\end{remark}

\begin{remark}
Railway alignments are embedded in coordinate systems, measurement
processes, construction workflows, maintenance histories, and
operational environments.
\end{remark}

\begin{remark}
Measurements are imperfect.

Coordinate systems may be ambiguous.

Constructed infrastructure deviates from design intent.

Observed geometry may differ from modeled geometry.
\end{remark}

\begin{remark}
The missing element is therefore the interaction between the ideal
alignment model and physical reality.
\end{remark}

% ============================================================
\section*{What comes next?}
% ============================================================

\begin{remark}
This part introduces the reality layer of AIM.
\end{remark}

\begin{remark}
The chapters that follow investigate:
\begin{itemize}
\item coordinate reference systems,
\item measurement and uncertainty,
\item stationing systems,
\item topology and special railway elements,
\item construction processes,
\item and realized railway geometry.
\end{itemize}
\end{remark}

\begin{remark}
Particular emphasis is placed on the distinction between
\emph{designed} and \emph{realized} railway states.
\end{remark}

\begin{remark}
This distinction ultimately leads to the concept of realization-aware
alignment modelling, one of the central ideas underlying AIM.
\end{remark}

% ============================================================
\section*{Relation to the AIM Roadmap}
% ============================================================

\begin{remark}
Within the AIM roadmap, the present part corresponds to the transition
from ideal runtime evaluation toward interaction with the physical
world.
\end{remark}

\begin{center}
\begin{tikzpicture}[
  node distance=1.4cm,
  box/.style={
    draw,
    rounded corners,
    minimum width=4.8cm,
    minimum height=0.9cm,
    align=center
  }
]

\node[box] (runtime)
{Runtime Services};

\node[box, below=of runtime] (reality)
{Reality Layer};

\node[box, below=of reality] (realization)
{Realization};

\draw[->, thick] (runtime) -- (reality);
\draw[->, thick] (reality) -- (realization);

\end{tikzpicture}
\end{center}

\begin{remark}
The progression
\[
\text{Runtime Services}
\rightarrow
\text{Reality Layer}
\rightarrow
\text{Realization}
\]
marks the transition from ideal state evaluation toward interaction
with measured and constructed railway systems.
\end{remark}

% ============================================================
\begin{principle}[Reality Principle]
\label{princ:reality-principle}
Within AIM, railway models and physical railway systems must be treated
as related but distinct entities connected through measurement,
construction, realization, and interpretation processes.
\end{principle}

% ============================================================
\begin{summary}
Geometry explains shape.

States explain configuration.

Runtime explains evaluation.

Reality explains how mathematical railway models interact with the
physical world.

The present part therefore establishes the bridge between ideal railway
states and realized railway infrastructure.
\end{summary}
